\section{Anillos japoneses}
\label{section:0.23-es}

Los resultados de esta sección se completarán en IV, (7.6) y (7.7).

\subsection{Anillos japoneses}
\label{subsection:0.23.1-es}

\begin{definition}[23.1.1]
\label{0.23.1.1-es}
Se dice que un anillo íntegro $A$ es un anillo japonés si, para toda
extensión finita $K'$ de su cuerpo de fracciones $K$, la clausura integral
$A'$ de $A$ en $K'$ es

\oldpage[0\textsubscript{IV-1}]{214}

un $A$-módulo de tipo finito (en otras palabras, una $A$-álgebra finita). Se dice
que un anillo $A$ es universalmente japonés si toda $A$-álgebra íntegra de
tipo finito es un anillo japonés.
\end{definition}

Está claro que si $A$ es un anillo japonés, todo anillo de fracciones
$S^{-1}A$ es entonces un anillo japonés, pues (con las notaciones precedentes)
$S^{-1}A'$ es la clausura integral de $S^{-1}A$ en $K'$, y es
evidentemente un $S^{-1}A$-módulo de tipo finito.

Un anillo íntegro noetheriano, incluso un \emph{anillo de valoración discreta},
no es necesariamente un anillo japonés [30].

\begin{proposition}[23.1.2]
\label{0.23.1.2-es}
Sean $A$ un anillo noetheriano íntegro y $K$ su cuerpo de fracciones. Si, para
toda extensión radicial finita $K'$ de $K$, la clausura integral de $A$
en $K'$ es un $A$-módulo de tipo finito, entonces $A$ es un anillo japonés.
\end{proposition}

\begin{proof}
Como $A$ es noetheriano, para comprobar que $A$ es un anillo japonés basta
ver que, para toda extensión cuasi-Galois finita $L$ de $K$, la clausura
integral $B$ de $A$ en $L$ es un $A$-módulo de tipo
finito. Ahora bien, $L$ es
una extensión galoisiana de la mayor extensión radicial $K'$ de $K$ contenida
en $L$; y si $A'$ es la clausura integral de $A$ en $K'$, $B$ es la clausura
integral de $A'$ en $L$. Pero se sabe que en una extensión separable de $K'$,
la clausura integral de $A'$ es un $A'$-módulo de tipo finito (Bourbaki,
\emph{Alg. comm.}, cap.~V, §~1, n\textsuperscript{o}~6, cor.~1 de la
prop.~20), de donde se sigue la proposición.
\end{proof}

De (23.1.2) se deduce que, cuando $K$ es de característica $0$, decir que
$A$ es un anillo japonés significa que su \emph{clausura integral} $A'$ es
un $A$-módulo de tipo finito.

\begin{theorem}[23.1.3]
\label{0.23.1.3-es}
\textnormal{(Tate).} Sean $A$ un anillo íntegro noetheriano y $x\neq0$ un
elemento de $A$. Supongamos que se cumplen las condiciones siguientes~:
\begin{enumerate}
  \item[(i)] $A$ es íntegramente cerrado.
  \item[(ii)] $\mathfrak p=xA$ es primo y $A$ es separado y completo para
  la topología $\mathfrak p$-preádica.
  \item[(iii)] $A/xA$ es un anillo japonés.
\end{enumerate}
Entonces $A$ es un anillo japonés.
\end{theorem}

Utilizaremos el lema siguiente~:

\begin{lemma}[23.1.3.1]
\label{0.23.1.3.1-es}
Sean $A$ un anillo y $x$ un elemento de $A$ no divisor de $0$ en $A$, tal
que $xA=\mathfrak p$ sea primo; entonces, para todo entero $n>0$, la imagen
inversa de $x^nA_{\mathfrak p}$ en $A$ es $x^nA$.
\end{lemma}

\begin{proof}
En efecto, supongamos que $b\in A$ es un elemento tal que
$b/1=x^na/s$ en $A_{\mathfrak p}$, donde $a\in A$ y
$s\notin\mathfrak p$; existe entonces $s'\notin\mathfrak p$ tal que
$s'sb=x^nas'$, de donde $s'sb\in\mathfrak p$, y como
$s's\notin\mathfrak p$, esto implica $b\in\mathfrak p$, es decir,
$b=xb'$ con $b'\in A$; puesto que $x$ no es divisor de $0$, se concluye
$s'sb'=x^{n-1}as$, y basta razonar por inducción sobre $n$.
\end{proof}

\begin{proof}
Para demostrar el teorema, podemos limitarnos al caso en que el cuerpo de fracciones
$K$ de $A$ es de característica $p>0$, puesto que $A$ es íntegramente cerrado.
Sea $K'$ una extensión radicial finita de $K$, de modo que existe una
potencia $q=p^e$ tal que $K'^q\subset K$; sustituyendo $K'$ por una
extensión radicial mayor, podemos incluso suponer que existe $y\in K'$ tal
que $y^q=x$. Sea $A'$ la clausura integral de $A$ en $K'$; como
$A$ es íntegramente cerrado, $A'$ es el conjunto de los $x'\in K'$ tales que
$x'^q\in A'\cap K=A$. Pongamos $V=A_{\mathfrak p}$; como el ideal maximal
$\mathfrak m=xV$ de $V$ es principal, se sabe que el anillo local
noetheriano $V$ es un \emph{anillo de valoración discreta} (17.1.4); como
$V$ es íntegramente cerrado, el mismo

\oldpage[0\textsubscript{IV-1}]{215}

razonamiento anterior muestra que la clausura integral $V'$ de $V$ en
$K'$ es el conjunto de los $x'\in K'$ tales que $x'^q\in V$; se sabe
(Bourbaki, \emph{Alg. comm.}, cap.~VI, §~8, n\textsuperscript{o}~6,
prop.~6, y cap.~V, §~2, n\textsuperscript{o}~3, lema~4) que $V'$ es un
\emph{anillo de valoración discreta}, cuyo ideal maximal $\mathfrak m'$ es
el conjunto de los $x'\in K'$ tales que $x'^q\in\mathfrak m$ y, además, el cuerpo
residual $V'/\mathfrak m'$ es una extensión \emph{finita} de
$V/\mathfrak m$ (\emph{loc. cit.}, cap.~VI, §~8,
n\textsuperscript{o}~1, lema~2). Demostremos que, para todo entero $n>0$, se tiene
\begin{equation}
\label{0.23.1.3.2-es}
  \mathfrak m'^n\cap A'=y^nA'.
  \tag{23.1.3.2}
\end{equation}

En efecto, como $y^q=x$, se tiene $y\in\mathfrak m'$, por tanto
$y^nA'\subset\mathfrak m'^n\cap A'$. Recíprocamente, sea
$x'\in\mathfrak m'^n\cap A'$, y pongamos $x'=z'y^n$ con $z'\in K'$. Se tiene
$x'^q\in A'$; por otra parte, $x'$ es suma de productos
$t'_1\cdots t'_n$ con $t'_i\in\mathfrak m'$, por tanto
$t_i'^q\in\mathfrak m$, y de ello se concluye que
$x'^q\in\mathfrak m^n\cap A$. Ahora bien, el lema (23.1.3.1) prueba que
$\mathfrak m^n\cap A=x^nA$, por lo que se tiene $z'^qx^n\in x^nA$, o también
$z'^q\in A$, puesto que $x\neq0$, lo que establece (23.1.3.2).

Demostremos, en segundo lugar, que sobre $A'$ la topología $xA'$-preádica es
\emph{separada}; esta topología es también la topología $yA'$-preádica de
$A'$, y de (23.1.3.2) se deduce que esta topología es inducida por la
topología $\mathfrak m'$-preádica de $V'$, que es separada porque $V'$ es
un anillo de valoración discreta.

Probemos ahora que $A'/xA'$ es un $A$-módulo de tipo finito. Como $x=y^q$
y $y^kA'/y^{k+1}A'$ es isomorfo a $A'/yA'$, podemos limitarnos a
demostrar que $A'/yA'$ es un $A$-módulo de tipo finito; pero la fórmula
(23.1.3.2), aplicada para $n=1$, muestra que $A'/yA'$ es un subanillo de
$V'/\mathfrak m'$, es decir, del cuerpo residual de $V'$, que es una
extensión finita del cuerpo residual $V/\mathfrak m$ de $V$. Ahora bien,
$V/\mathfrak m$ es el cuerpo de fracciones de $A/\mathfrak p$, y como
$A'$ es entero sobre $A$, $A'/yA'$ es entero sobre $A/\mathfrak p$, por tanto
está contenido en la clausura integral de $A/\mathfrak p$ en
$V'/\mathfrak m'$; como por hipótesis $A/\mathfrak p$ es un anillo
japonés noetheriano, $A'/yA'$ es un $(A/\mathfrak p)$-módulo de tipo finito,
y \emph{a fortiori} un $A$-módulo de tipo finito.

Así pues, puesto que la topología $xA'$-preádica de $A'$ es separada, $A'$ es un
subanillo de su completado $\widehat{A'}$ para esta topología; pero, puesto que
$A'/xA'$ es un $A$-módulo de tipo finito y $A$ es separado y completo para
la topología $xA$-preádica, $\widehat{A'}$ es un $A$-módulo de tipo finito en
virtud de (0\textsubscript{I}, 7.2.9), por tanto lo mismo vale para $A'$, C.Q.D.
\end{proof}

\begin{corollary}[23.1.4]
\label{0.23.1.4-es}
Sea $A$ un anillo japonés noetheriano e íntegramente cerrado. Entonces todo
anillo de series formales $A[[T_1,\ldots,T_r]]$ es un anillo japonés.
\end{corollary}

\begin{proof}
Se sabe que, para todo anillo noetheriano íntegramente cerrado $B$, el anillo de
series formales $B[[T]]$ es noetheriano e íntegramente cerrado (Bourbaki,
\emph{Alg. comm.}, cap.~V, §~1, n\textsuperscript{o}~4, prop.~14); por
tanto, por inducción sobre $n$, podemos limitarnos a probar que $A[[T]]$ es un anillo
japonés. Ahora bien, el elemento $x=T$ cumple todas las condiciones de (23.1.3), de donde
se sigue la conclusión.
\end{proof}

\begin{theorem}[23.1.5]
\label{0.23.1.5-es}
\textnormal{(Nagata).} Todo anillo local noetheriano íntegro completo $A$ es
un anillo japonés.
\end{theorem}

\begin{proof}
Se sabe (19.8.8, (ii)) que existe un subanillo $B$ de $A$ que es un
anillo local completo y regular, tal que $A$ es una $B$-álgebra finita; basta
evidentemente probar

\oldpage[0\textsubscript{IV-1}]{216}

que $B$ es un anillo japonés; en otras palabras, podemos limitarnos al caso en que
$A$ es además \emph{regular}. Razonemos por inducción sobre
$n=\dim(A)$; el teorema es evidente para $n=0$. Supongamos, pues, $n>0$ y,
denotando por $\mathfrak m$ el ideal maximal de $A$, sea $x$ un elemento de
$\mathfrak m-\mathfrak m^2$; se sabe (17.1.8) que $A/xA$ es un anillo
regular, y además ((17.1.7) y (16.3.4)) que
$\dim(A/xA)=n-1$; además, el ideal $xA$ es cerrado en $A$
(0\textsubscript{I}, 7.3.5), por lo que $A/xA$ es completo. En virtud de la hipótesis
de inducción, $A/xA$ es, por tanto, un anillo japonés, y de
(23.1.3) se deduce que lo mismo vale para $A$.
\end{proof}

\begin{corollary}[23.1.6]
\label{0.23.1.6-es}
Sean $A$ un anillo local noetheriano íntegro completo y $K$ su cuerpo de
fracciones; sea $K'$ una extensión finita de $K$; entonces la clausura integral
$A'$ de $A$ en $K'$ es un anillo local completo.
\end{corollary}

\begin{proof}
Ya se sabe por (23.1.5) que $A'$ es una $A$-álgebra finita, por tanto un anillo
semilocal noetheriano completo y, en consecuencia, un \emph{producto directo}
de anillos locales; pero como $A'$ es íntegro, es necesariamente un
anillo local.
\end{proof}

\begin{proposition}[23.1.7]
\label{0.23.1.7-es}
Sean $A$ un anillo local noetheriano íntegro y $K$ su cuerpo de fracciones;
sean $K'$ una extensión finita de $K$ y $A'$ la clausura integral de $A$ en
$K'$. Se supone que el completado $\widehat A$ es reducido, y se denota por
$R$ su anillo total de fracciones.
\begin{enumerate}
  \item[(i)] Si el anillo $R\otimes_KK'$ es reducido (lo que ocurrirá en
  particular cuando $K'$ sea una extensión separable de $K$, siendo $R$ un
  producto directo de cuerpos que son extensiones de $K$ (Bourbaki, \emph{Alg.},
  cap.~VIII, §~7, n\textsuperscript{o}~3, cor.~1 del teo.~1)), entonces $A'$ es
  un $A$-módulo de tipo finito.
  \item[(ii)] Si, en particular, $R$ es una $K$-álgebra separable, entonces $A$
  es un anillo japonés.
\end{enumerate}
\end{proposition}

\begin{proof}
\emph{(i)} Pongamos, para simplificar, $A_1=\widehat A$,
$A'_1=A'\otimes_AA_1$, $K'_1=K'\otimes_AA_1$. Como $A_1$ es un
$A$-módulo plano (0\textsubscript{I}, 7.3.5), $A'_1$ se identifica con un
subanillo de $K'_1$ y es evidentemente entero sobre $A_1$. Por otra parte, si
se pone $K_1=K\otimes_AA_1$, se puede escribir
$K'_1=K'\otimes_KK_1$; como $A$ es íntegro y $A_1$ es un $A$-módulo
plano, todo elemento $\neq0$ de $A$ es $A_1$-regular
(0\textsubscript{I}, 6.3.4); como $K=S^{-1}A$, donde $S=A-\{0\}$, se tiene
$K_1=S^{-1}A_1$, y la observación anterior prueba que $K_1$ se identifica con un
subanillo del anillo total de fracciones $R$ de $A_1$. Puesto que todo
$K$-módulo es plano, $K'_1$ se identifica con un subanillo de $K'\otimes_KR$,
que por hipótesis es reducido y es una $R$-álgebra finita. Designemos por
$\mathfrak q_i$ ($1\leqslant i\leqslant n$) los ideales primos minimales de
$A_1$, por $B_i$ el anillo íntegro $A_1/\mathfrak q_i$, por $L_i$ el cuerpo de
fracciones de $B_i$, de modo que $R$ es producto directo de los $L_i$, y
$K'\otimes_KR$ producto directo de los $K'\otimes_KL_i$; estas últimas
$K$-álgebras, al ser reducidas, son productos directos de cuerpos,
extensiones finitas de los $L_i$. Como $B_i$ es un anillo local íntegro
\emph{completo}, el teorema de Nagata (23.1.5) prueba que la clausura
integral de $B_i$ en $K'\otimes_KL_i$ es un $B_i$-módulo de tipo finito, y
\emph{a fortiori} un $A_1$-módulo de tipo finito; como todo elemento de
$K'\otimes_KL_i$ que es entero sobre $A_1$ es también entero sobre $B_i$ (al ser
anulado por $\mathfrak q_i$), se concluye finalmente que la clausura
integral de $A_1$ en $K'\otimes_KR$ es un $A_1$-módulo de tipo finito. Pero
como $A'_1$ está contenido en esta clausura integral y $A_1$ es
noetheriano, $A'_1$ es también un $A_1$-módulo de tipo finito. Por último, como $A_1$
es un $A$-módulo fielmente plano (0\textsubscript{I}, 7.3.5), se concluye
que $A'$ es un $A$-módulo de tipo finito (Bourbaki, \emph{Alg. comm.},
cap.~I, §~3, n\textsuperscript{o}~6, prop.~11).

\oldpage[0\textsubscript{IV-1}]{217}

\emph{(ii)} La hipótesis implica que los $L_i$ son extensiones separables
de $K$, por tanto $K'\otimes_KL_i$ es reducido (Bourbaki, \emph{Alg.}, cap.~VIII,
§~7, n\textsuperscript{o}~3, teo.~1) y se puede aplicar (i) a toda
extensión finita $K'$ de $K$, lo que demuestra nuestra afirmación.
\end{proof}

\subsection{Clausura integral de un anillo local noetheriano íntegro}
\label{subsection:0.23.2-es}

\begin{env}[23.2.1]
\label{0.23.2.1-es}
En lo que sigue, para todo anillo $A$, denotaremos por $A_{\mathrm{red}}$, para
abreviar, el cociente de $A$ por su nilradical (de modo que si
$X=\operatorname{Spec}(A)$, se tiene
$X_{\mathrm{red}}=\operatorname{Spec}(A_{\mathrm{red}})$).

Diremos que un anillo \emph{local} $A$ es \emph{unirrama} si el anillo
$A_{\mathrm{red}}$ es íntegro y si la clausura integral de
$A_{\mathrm{red}}$ es un anillo \emph{local}, lo que generaliza la definición
dada en (III, 4.3.6). Diremos que $A$ es \emph{geométricamente
unirrama} si es unirrama y si el cuerpo residual del anillo local,
clausura integral de $A_{\mathrm{red}}$, es una extensión \emph{radicial}
del de $A$. Es claro que un anillo local íntegramente cerrado es
geométricamente unirrama~; de (23.1.6) resulta que un anillo local
noetheriano íntegro \emph{completo} es unirrama.
\end{env}

\begin{lemma}[23.2.2]
\label{0.23.2.2-es}
\textnormal{(*)}\footnote{En lo que sigue de este número, utilizamos
nociones y resultados expuestos en el cap.~IV, §§~2 y 5~; como los resultados de
este número no se utilizan antes del cap.~IV, §~6, esto no da lugar a un
círculo vicioso.} Sean $A$ un anillo íntegro, $A'$ su clausura integral~; para
que el morfismo canónico
$f:\operatorname{Spec}(A')\to\operatorname{Spec}(A)$ sea radicial, es necesario y
suficiente que para todo $\mathfrak p\in\operatorname{Spec}(A)$,
$A_{\mathfrak p}$ sea geométricamente unirrama~; cuando así ocurre,
$f$ es un homeomorfismo.
\end{lemma}

\begin{proof}
En efecto, para todo $\mathfrak p\in\operatorname{Spec}(A)$, la clausura
integral de $A_{\mathfrak p}$ es $A'_{\mathfrak p}$ (Bourbaki,
\emph{Alg. comm.}, cap.~V, §~1, n\textsuperscript{o}~5, prop.~16), y todos
los ideales primos de $A'_{\mathfrak p}$ situados sobre
$\mathfrak pA_{\mathfrak p}$ son maximales (\emph{loc. cit.}, §~2,
n\textsuperscript{o}~1, prop.~1). Decir que $f$ es inyectivo significa, por tanto, que
para todo $\mathfrak p\in\operatorname{Spec}(A)$, $A'_{\mathfrak p}$ es un
anillo \emph{local}, es decir, que cada uno de los $A_{\mathfrak p}$ es
unirrama. Decir que todo $A_{\mathfrak p}$ es geométricamente unirrama
significa entonces que $f$ es radicial en virtud de (I, 3.5.8). Cuando así
ocurre, $f$ es sobreyectivo y cerrado (II, 6.1.10), luego es un homeomorfismo.
\end{proof}

\begin{lemma}[23.2.3]
\label{0.23.2.3-es}
Sean $A$ un anillo íntegro, $K$ su cuerpo de fracciones, $B$ un anillo
noetheriano, $\varphi:A\to B$ un homomorfismo de anillos que hace de $B$ un
$A$-módulo plano. Sea $K'$ una extensión de $K$~; pongamos
$B_1=B_{\mathrm{red}}$, y sea $R$ el anillo total de fracciones de $B_1$.
Entonces, para toda sub-$A$-álgebra $A'$ de $K'$, el homomorfismo canónico
\begin{equation}
\label{0.23.1.8.1-es}
  (A'\otimes_AB)_{\mathrm{red}}
  \longrightarrow(K'\otimes_AR)_{\mathrm{red}}
  \tag{23.1.8.1}
\end{equation}
es inyectivo.
\end{lemma}

\begin{proof}
En efecto, se puede considerar el homomorfismo canónico
$h:A'\otimes_AB\to K'\otimes_AR$ como compuesto de los homomorfismos siguientes
\[
  A'\otimes_AB\xrightarrow{u}K'\otimes_AB
  \xrightarrow{v}K'\otimes_AB_1\xrightarrow{w}K'\otimes_AR.
\]
Como $B$ es un $A$-módulo plano, $u$ es inyectivo~; asimismo, como $K'$ es un
$K$-módulo plano y $K$ un $A$-módulo plano, $K'$ es un $A$-módulo plano, y
$w$ es, por tanto, inyectivo, puesto que $B_1$ es reducido y, por consiguiente, el homomorfismo
$B_1\to R$ es inyectivo. Por último, por la misma razón, si $\mathfrak N$ es

\oldpage[0\textsubscript{IV-1}]{218}

el nilradical de $B$, el núcleo de $v$ es $K'\otimes_A\mathfrak N$, por lo que está
contenido en el nilradical de $K'\otimes_AB$~; se concluye inmediatamente que si
la imagen por $h=w\circ v\circ u$ de un elemento $x\in A'\otimes_AB$ es
nilpotente, $x$ es nilpotente, lo que demuestra el lema.
\end{proof}

\begin{proposition}[23.2.4]
\label{0.23.2.4-es}
Sean $A$ un anillo íntegro, $K$ su cuerpo de fracciones, $B$ un anillo
noetheriano, $\mathfrak q_i$ ($1\leqslant i\leqslant n$) sus ideales primos
minimales, $\varphi:A\to B$ un homomorfismo de anillos. Se supone que los
$B/\mathfrak q_i$ son anillos japoneses y que $B$ es un $A$-módulo
plano. Sea $K'$ una extensión finita de $K$, y sea $(C_\lambda)$ la familia
filtrante creciente de subanillos de $K'$ que son $A$-álgebras
finitas y admiten $K'$ como cuerpo de fracciones~; la unión $A'$ de los
$C_\lambda$ es, por tanto, la clausura integral de $A$ en $K'$. Entonces~:
\begin{enumerate}
  \item[(i)] Existe un índice $\alpha$ tal que para $\lambda\geqslant
  \alpha$, el homomorfismo canónico
  \begin{equation}
  \label{0.23.2.4.1-es}
    (C_\alpha\otimes_AB)_{\mathrm{red}}
    \longrightarrow(C_\lambda\otimes_AB)_{\mathrm{red}}
    \tag{23.2.4.1}
  \end{equation}
  es biyectivo.
  \item[(ii)] Si además $B$ es un $A$-módulo fielmente plano, el morfismo
  $\operatorname{Spec}(A')\to\operatorname{Spec}(C_\alpha)$ es radicial.
\end{enumerate}
\end{proposition}

\begin{proof}
\emph{(i)} Sea $B_1=B_{\mathrm{red}}$, y sea $R$ el anillo total de
fracciones de $B_1$~; como $A$ es íntegro y $B$ un $A$-módulo plano, todo
elemento $\neq0$ de $A$ es $B$-regular (0\textsubscript{I}, 6.3.4), por lo que
el homomorfismo compuesto $A\xrightarrow{\varphi}B\to B_1$ se prolonga a un
homomorfismo $K\to R$~; se puede escribir entonces
$K'\otimes_AR=(K'\otimes_AK)\otimes_KR$, y como $K'\otimes_AK$
se identifica con $K'$, se tiene $K'\otimes_AR=K'\otimes_KR$. Como $R$ es producto
directo de los cuerpos de fracciones $L_i$ de los $B/\mathfrak q_i$,
$K'\otimes_AR$ es producto directo de los $K'\otimes_KL_i$, y por consiguiente
$(K'\otimes_AR)_{\mathrm{red}}$ es producto directo de un número finito
de extensiones finitas de los $L_i$. En virtud de la hipótesis, la clausura
integral de $B/\mathfrak q_i$ en una extensión finita de $L_i$ es un
$(B/\mathfrak q_i)$-módulo de tipo finito, luego un $B$-módulo de tipo finito~;
se concluye que la clausura integral $B'$ de $B$ en
$(K'\otimes_AR)_{\mathrm{red}}$ es un $B$-módulo de tipo finito. En virtud de
(23.2.3), los $(C_\lambda\otimes_AB)_{\mathrm{red}}$ se identifican
canónicamente con subanillos de $(K'\otimes_AR)_{\mathrm{red}}$ que son
$B$-álgebras finitas, y están, por tanto, contenidos en $B'$. Como $B$ es noetheriano
y $B'$ un $B$-módulo de tipo finito, la familia filtrante de los
$(C_\lambda\otimes_AB)_{\mathrm{red}}$ admite un \emph{elemento mayor}
$(C_\alpha\otimes_AB)_{\mathrm{red}}$, lo que demuestra (i).

\emph{(ii)} Como $B$ es un $A$-módulo fielmente plano, basta
(IV, 2.6.1, (v)) demostrar que el morfismo
$\operatorname{Spec}(A'\otimes_AB)\to
\operatorname{Spec}(C_\alpha\otimes_AB)$ es radicial, o, lo que es equivalente
(I, 5.1.6), que el morfismo
$\operatorname{Spec}(A'\otimes_AB)_{\mathrm{red}}\to
\operatorname{Spec}(C_\alpha\otimes_AB)_{\mathrm{red}}$ es radicial~; pero
se tiene $A'\otimes_AB=\varinjlim_\lambda(C_\lambda\otimes_AB)$, luego
(IV, 5.13.2)
$(A'\otimes_AB)_{\mathrm{red}}=
\varinjlim_\lambda(C_\lambda\otimes_AB)_{\mathrm{red}}$, y la conclusión
resulta de (i).
\end{proof}

\begin{corollary}[23.2.5]
\label{0.23.2.5-es}
Sean $A$ un anillo local noetheriano íntegro, $K$ su cuerpo de fracciones,
$K'$ una extensión finita de $K$. Sea $(C_\lambda)$ la familia filtrante
creciente de subanillos de $K'$ que son $A$-álgebras finitas (por tanto,
anillos noetherianos semilocales (Bourbaki, \emph{Alg. comm.}, cap.~IV,
§~2, n\textsuperscript{o}~5, cor.~3 de la prop.~9)) y admiten $K'$ como
cuerpo de fracciones. Entonces existe $\alpha$ tal que el homomorfismo
$(\widehat C_\alpha)_{\mathrm{red}}\to
(\widehat C_\lambda)_{\mathrm{red}}$ sea un isomorfismo para
$\lambda\geqslant\alpha$, y si $A'$ es la clausura integral de $A$ en
$K'$, el morfismo
$\operatorname{Spec}(A')\to\operatorname{Spec}(C_\alpha)$ es radicial
(cf. (23.2.2)).
\end{corollary}

\oldpage[0\textsubscript{IV-1}]{219}

\begin{proof}
Se aplica (23.2.4) tomando $B=\widehat A$~; como los
$B/\mathfrak q_i$ son anillos locales noetherianos completos, son anillos
japoneses (23.1.5) y $B$ es un $A$-módulo fielmente plano
(0\textsubscript{I}, 7.3.5)~; además, se tiene
$C_\lambda\otimes_AB=\widehat C_\lambda$ (Bourbaki, \emph{Alg. comm.},
cap.~III, §~3, n\textsuperscript{o}~4, th.~3 y cap.~IV, §~2,
n\textsuperscript{o}~5, cor.~3 de la prop.~9).
\end{proof}

\begin{corollary}[23.2.6]
\label{0.23.2.6-es}
Bajo las hipótesis de (23.2.5), la clausura integral $A'$ de $A$ en
$K'$ es un anillo semilocal~; si $\mathfrak m$ es el ideal maximal de $A$,
entonces, para todo ideal maximal $\mathfrak m'$ de $A'$,
$A'/\mathfrak m'$ es una extensión finita de $A/\mathfrak m$.
\end{corollary}

\begin{proof}
El hecho de que el homomorfismo
$(\widehat C_\alpha)_{\mathrm{red}}\to
(\widehat C_\lambda)_{\mathrm{red}}$ sea biyectivo para
$\lambda\geqslant\alpha$ implica que el número de ideales maximales de
$C_\lambda$ es constante para $\lambda\geqslant\alpha$ y que si
$\mathfrak m_\alpha$ es un ideal maximal de $C_\alpha$ y
$\mathfrak m_\lambda$ el único ideal maximal de $C_\lambda$ situado sobre
$\mathfrak m_\alpha$, los cuerpos $C_\alpha/\mathfrak m_\alpha$ y
$C_\lambda/\mathfrak m_\lambda$ son canónicamente isomorfos. La conclusión
resulta de que $\mathfrak m'=\varinjlim\mathfrak m_\lambda$ y
$A'/\mathfrak m'=\varinjlim(C_\lambda/\mathfrak m_\lambda)$
(0\textsubscript{III}, 10.3.1.1.3).
\end{proof}

\begin{proposition}[23.2.7]
\label{0.23.2.7-es}
\textnormal{(Y. Mori).} Sean $A$ un anillo local noetheriano íntegro, $K$ su
cuerpo de fracciones, $A'$ la clausura integral de $A$. Entonces $A'$ es un
anillo de Krull (Bourbaki, \emph{Alg. comm.}, cap.~VII, §~1) semilocal~;
en otras palabras, existe una familia $(V_\lambda)$ de anillos de valoración
discreta que tienen $K$ como cuerpo de fracciones, tales que
$A'=\bigcap_\lambda V_\lambda$ y que todo $x\in K$ pertenece a todos los
$V_\lambda$ salvo a un número finito de ellos. Además, existe un
subanillo $C$ de $K$ que es un $A$-álgebra finita, tal que los
$V_\lambda$ son las clausuras integrales de los anillos
$C_{\mathfrak p_\lambda}$, donde $(\mathfrak p_\lambda)$ es la familia de los
ideales primos de altura $1$ del anillo local $C$.
\end{proposition}

Demostremos primero el lema siguiente~:

\begin{lemma}[23.2.7.1]
\label{0.23.2.7.1-es}
Sean $A$, $B$ dos anillos, $\varphi:A\to B$ un homomorfismo que hace de $B$
un $A$-módulo fielmente plano. Se supone que $A$ es íntegro~; entonces, si
$C=B_{\mathrm{red}}$, el homomorfismo compuesto
$\psi:A\to B\to B_{\mathrm{red}}=C$ es inyectivo y se prolonga a un
homomorfismo inyectivo del cuerpo de fracciones $K$ de $A$ en el anillo total
de fracciones $R$ de $C$~; además, si $C'$ es la clausura integral de
$C$ en $R$, $C'\cap K$ es la clausura integral de $A$.
\end{lemma}

\begin{proof}
Como $A$ es íntegro y $\varphi$ es inyectivo, la intersección de $\varphi(A)$
y del nilradical $\mathfrak N$ de $B$ se reduce a $0$, y como
$C=B/\mathfrak N$, $\psi$ es inyectivo. Todo $a\neq0$ de $A$, al no ser
divisor de cero en $A$, tampoco lo es en $B$ por planitud
(0\textsubscript{I}, 6.3.4)~; se deduce que $a$ tampoco es divisor
de $0$ en $C$, pues si se tuviera $ax\in\mathfrak N$ para un
$x\notin\mathfrak N$ de $B$, se obtendría $a^nx^n=0$ para un entero $n$,
lo que contradice lo anterior, puesto que $x^n\neq0$. Se puede, por tanto, prolongar
$\psi$ a un homomorfismo inyectivo de $K$ en $R$. Para demostrar la última
afirmación, observemos que es claro que la clausura integral $A'$ de $A$ está
contenida en $C'\cap K$. Recíprocamente, sea $x\in C'\cap K$~; $x$ es
entero sobre $C$ y, \emph{a fortiori}, sobre $B$, es decir, $B[x]$ es una
$B$-álgebra finita~; además, $K$ se identifica con un subanillo de
$B\otimes_AK$, y el subanillo $B[x]$ de $B\otimes_AK$ se identifica con
$B\otimes_AA[x]$ por planitud~; se concluye que $A[x]$ es un $A$-módulo
de tipo finito (Bourbaki, \emph{Alg. comm.}, cap.~I\textsuperscript{er}, §~3,
n\textsuperscript{o}~6, prop.~11), y, por tanto, que $x\in A'$.
\end{proof}

Establecido este lema, vamos a aplicarlo a la inyección canónica de $A$
en su completado $\widehat A$, que es un $A$-módulo fielmente plano~;
$B=(\widehat A)_{\mathrm{red}}=\widehat A/\mathfrak N$ (donde $\mathfrak N$ es
el nilradical de $\widehat A$) es un anillo local noetheriano completo reducido,
cuyo anillo total de

\oldpage[0\textsubscript{IV-1}]{220}

fracciones $R$ es, por tanto, producto directo de un número finito de cuerpos $L_i$~;
la imagen canónica $B_i$ de $B$ en $L_i$ es un anillo local noetheriano
\emph{íntegro y completo}, cuyo cuerpo de fracciones es $L_i$, y la
clausura integral $B'$ de $B$ en $R$ es \emph{producto directo} de los
$B'_i$, donde $B'_i$ es la clausura integral de $B_i$. Pero en virtud de
(23.1.5), $B'_i$ es una $B_i$-álgebra \emph{finita}, y, por tanto, un anillo local
\emph{noetheriano} íntegramente cerrado. Se sabe (Bourbaki,
\emph{Alg. comm.}, cap.~VII, §~1, n\textsuperscript{o}~3, cor. del th.~2) que
$B'_i$ es un anillo de Krull. Ahora bien, para todo $i$, se tiene un homomorfismo
$\varphi_i:K\to L_i$, que es inyectivo, y por consiguiente
$\varphi_i^{-1}(B'_i)$ es un anillo de Krull en $K$. Pero como
$A'=B'\cap K$ en virtud del lema y
$B'=\bigcap_i\varphi_i^{-1}(B'_i)$, $A'$ es intersección de un número finito
de anillos de Krull y, por tanto, es un anillo de Krull.

Para demostrar la última afirmación de la proposición, observemos que existe una
$A$-álgebra finita $C\subset K$ tal que el morfismo
$\operatorname{Spec}(A')\to\operatorname{Spec}(C)$ es radicial y un
homeomorfismo (23.2.4)~; para todo ideal primo
$\mathfrak p'\in\operatorname{Spec}(A')$, $\mathfrak p'$ es, por tanto, el único
ideal primo de $A'$ situado sobre $\mathfrak p=\mathfrak p'\cap C$ y
$A'_{\mathfrak p'}$ es la clausura integral de $C_{\mathfrak p}$~; por otra parte,
el hecho de que $\operatorname{Spec}(A')\to\operatorname{Spec}(C)$ sea un
homeomorfismo implica que la aplicación
$\mathfrak p'\mapsto\mathfrak p'\cap C$ es una biyección del conjunto de los
ideales primos de altura $1$ de $A'$ sobre el conjunto de los ideales primos de
altura $1$ de $C$~; de donde resulta la conclusión, teniendo en cuenta Bourbaki,
\emph{Alg. comm.}, cap.~VII, §~1, n\textsuperscript{o}~6, th.~4.

\begin{remarks}[23.2.8]
\label{0.23.2.8-es}
\begin{enumerate}
  \item[(i)] Conviene observar, al aplicar (23.2.6) y
  (23.2.7), que el anillo $A'$ \emph{no es necesariamente un anillo
  noetheriano}, como muestra un ejemplo de Nagata con $\dim(A)=3$ [30].

  \item[(ii)] Las conclusiones de (23.2.7) siguen siendo válidas si $A'$ es
  la clausura integral de $A$ en una extensión finita $K'$ de $K$~; en efecto,
  basta considerar una $A$-álgebra finita $B$ cuyo cuerpo de fracciones sea
  $K'$~; $B$ es un anillo semilocal noetheriano, y, por tanto,
  intersección de un número finito de anillos locales noetherianos
  $B_{\mathfrak m_i}$ ($\mathfrak m_i$ ideales maximales de $B$), y su clausura
  integral (que es igual a $A'$) es intersección de los
  $A'_{\mathfrak m_i}$ (Bourbaki, \emph{Alg. comm.}, cap.~II, §~3,
  n\textsuperscript{o}~3, cor.~4 del th.~1), que son las clausuras integrales
  de los $B_{\mathfrak m_i}$, y, por tanto, anillos de Krull por (23.1.13)~; por
  consiguiente, $A'$ es un anillo de Krull.

  \item[(iii)] Se puede demostrar ([30], 3.3.10) que para todo anillo noetheriano
  íntegro $A$ (no necesariamente local), la clausura integral de $A$ es un
  anillo de Krull.
\end{enumerate}
\end{remarks}

\begin{corollary}[23.2.9]
\label{0.23.2.9-es}
Sean $A$ un anillo noetheriano íntegro, $A'$ su clausura integral. Supongamos
que para todo anillo $C$ tal que $A\subset C\subset A'$ y tal que $C$ sea
una $A$-álgebra finita, y para todo ideal primo $\mathfrak q$ de altura $1$
en $C$, $\mathfrak q\cap A$ sea un ideal primo de altura $1$ en $A$.
Sea $P$ el conjunto de los ideales primos de altura $1$ de $A$~; se tiene entonces
\[
  A'=\bigcap_{\mathfrak p\in P}A'_{\mathfrak p}.
\]
\end{corollary}

\begin{proof}
Como $A'$ es un $A$-módulo sin torsión, se tiene
$A'=\bigcap A'_{\mathfrak m}$, donde $\mathfrak m$ recorre el conjunto de los ideales
maximales de $A$ (Bourbaki, \emph{Alg. comm.}, cap.~II, §~3,
n\textsuperscript{o}~3, cor.~4 del th.~1)~;

\oldpage[0\textsubscript{IV-1}]{221}

bastará, por tanto, demostrar que $A'_{\mathfrak m}$ es la intersección de los
$A'_{\mathfrak p}$ para los ideales primos $\mathfrak p\subset\mathfrak m$
de altura $1$. Como $A'_{\mathfrak m}$ es la clausura integral de
$A_{\mathfrak m}$, se ve que podemos limitarnos a demostrar el corolario para
$A_{\mathfrak m}$. Ahora bien, en virtud de (23.2.7), existe una
$A_{\mathfrak m}$-álgebra finita $B$ contenida en $A'_{\mathfrak m}$, tal
que $A'_{\mathfrak m}$ es la intersección de las clausuras integrales de los anillos
$B_{\mathfrak q}$, donde $\mathfrak q$ recorre el conjunto de los ideales primos de
altura $1$ de $B$. Pero $B$ es de la forma $C_{\mathfrak m}$, donde $C$ es
una $A$-álgebra finita~; por tanto, para todo ideal primo $\mathfrak q$ de altura
$1$ en $B$, $\mathfrak q\cap A$ es, por hipótesis, un ideal primo de
altura $1$ en $A$~; como la clausura integral de $B_{\mathfrak q}$
contiene la de $A_{\mathfrak q\cap A}$ y esta última contiene
$A'_{\mathfrak m}$, resulta que $A'_{\mathfrak m}$ es efectivamente
la intersección de los $A'_{\mathfrak p}$ para $\mathfrak p\in P$ y
$\mathfrak p\subset\mathfrak m$, lo que termina la demostración.
\end{proof}

\begin{remarks}[23.2.10]
\label{0.23.2.10-es}
\begin{enumerate}
  \item[(i)] Veremos (IV, 5.6.3 y 5.6.10) que la hipótesis de (23.2.9)
  se cumple cuando $A$ es \emph{universalmente catenario}, en particular
  (IV, 5.10.17 y 5.11.2) cuando $A$ es \emph{cociente de un anillo
  regular}~; por último, se cumple evidentemente si
  $\operatorname{Spec}(A')\to\operatorname{Spec}(A)$ es un homeomorfismo,
  en particular si este morfismo es radicial (23.2.2).

  \item[(ii)] La conclusión de (23.2.9) ya no es necesariamente cierta
  cuando se omite la hipótesis sobre las sub-$A$-álgebras finitas de $A'$. En
  efecto, se puede definir un anillo local noetheriano íntegro $A$, de dimensión
  $2$, cuya clausura integral $A'$ es una $A$-álgebra finita, pero tal que
  existe un ideal primo $\mathfrak p'$ de $A'$ de altura $1$ situado sobre
  el ideal maximal $\mathfrak m$ (de altura $2$) de $A$, y tal además que
  para todo ideal primo $\mathfrak p$ de altura $1$ de $A$,
  $A_{\mathfrak p}$ es íntegramente cerrado (IV, 5.6.11)~; la intersección de
  estos anillos $A_{\mathfrak p}$ no puede entonces ser igual a $A'$, puesto que $A'$
  es un anillo de Krull (Bourbaki, \emph{Alg. comm.}, cap.~VII, §~1,
  n\textsuperscript{o}~5, prop.~9).
\end{enumerate}
\end{remarks}

\begin{flushright}
\emph{(Continuará.)}
\end{flushright}
